% ================================================================
%  CHAPTER 1: INTRODUCTION
% ================================================================
\chapter{Introduction}

% ----------------------------------------------------------------
\section{Background}
% ----------------------------------------------------------------

Cancer remains one of the most formidable public health challenges of the
twenty-first century, imposing an extraordinary clinical, economic, and societal
burden on populations worldwide. According to the GLOBOCAN 2022 estimates compiled
by the International Agency for Research on Cancer (IARC), there were approximately
20 million new cancer cases and 9.7 million cancer-related deaths globally in 2022
\parencite{siegel2023cancer}. Among all cancer subtypes, \textbf{lung cancer} is
the leading cause of cancer mortality, accounting for approximately 1.8 million
deaths and representing nearly 18\% of all cancer fatalities in that year. It is
the second most commonly diagnosed cancer by incidence and the most lethal, primarily
because the majority of patients are diagnosed at advanced or metastatic stages when
curative surgery is no longer feasible.

\textbf{Non-Small Cell Lung Cancer (NSCLC)} is the predominant histological
category of lung cancer, constituting approximately 85\% of all diagnoses.
NSCLC encompasses three major subtypes: adenocarcinoma ($\sim$40\%), squamous cell
carcinoma ($\sim$25--30\%), and large-cell carcinoma ($\sim$10--15\%). The remaining
15\% is classified as Small Cell Lung Cancer (SCLC), a neuroendocrine tumour with
distinct biology and treatment approaches. The epidemiology of NSCLC has important
geographical and ethnic dimensions: the disease burden is particularly acute across
South and East Asia, including Nepal, where limited access to computed tomography
(CT) screening programmes and low public awareness result in late-stage diagnosis
in the majority of patients, precluding surgical intervention and significantly
worsening prognosis \parencite{rotow2023egfr}. The five-year overall survival rate
for stage IV NSCLC remains below 10\% with conventional chemotherapy.

\subsection{The Epidermal Growth Factor Receptor in NSCLC}

The \textbf{Epidermal Growth Factor Receptor (EGFR)} is a transmembrane receptor tyrosine kinase (RTK). Upon ligand binding, EGFR undergoes dimerization and intrinsic kinase activation, which initiates key intracellular downstream pathways—specifically \mbox{RAS–RAF–MEK–ERK}, \mbox{PI3K–AKT–mTOR}, STAT3/5, and \mbox{PLC-$\gamma$}—that orchestrate cellular proliferation, motility, and survival \parencite{rotow2023egfr}. In non-small cell lung cancer (NSCLC), somatic activating mutations within the EGFR kinase domain induce constitutive, ligand-independent receptor activation, thereby acting as primary oncogenic drivers. Two major sensitizing mutations account for 85–90\% of all EGFR-mutant NSCLC cases: \textbf{exon 19 deletions} (Ex19del; $\sim$45–50\%) and the \textbf{L858R} point mutation in the activation loop ($\sim$40\%).

EGFR-mutation prevalence varies sharply by ethnicity---10--15\% in Caucasian NSCLC
versus 40--50\% in East Asian populations \parencite{leonetti2019resistance}. Given
Nepal's rising NSCLC incidence and its admixed South/East Asian population, a large
and growing share of Nepali patients harbour targetable EGFR mutations, yet molecular
diagnostics and novel targeted therapies remain scarce---making affordable
computational drug discovery particularly relevant.

\subsection{Generations of EGFR Tyrosine Kinase Inhibitors}

EGFR TKIs have developed in three generations, each overcoming resistance to its
predecessor. \textbf{First-generation} reversible inhibitors (gefitinib, erlotinib)
achieved 55--75\% response rates and 9--13-month progression-free survival in
sensitising-mutant NSCLC, but resistance emerged in nearly all patients---most often
($\sim$50--60\%) via the \textbf{T790M} gatekeeper mutation.
\textbf{Second-generation} covalent inhibitors (afatinib, dacomitinib) bind Cys797
irreversibly but their wild-type EGFR inhibition causes dose-limiting toxicity.
\textbf{Third-generation} covalent inhibitors, led by \textbf{osimertinib}
(FDA-approved 2015), selectively target T790M-mutant over wild-type EGFR and, with
superior efficacy (PFS 18.9 vs.\ 10.2 months) and CNS penetration, are now the
standard of care \parencite{rotow2023egfr}.

\textbf{The C797S resistance problem.} Despite osimertinib's superiority, acquired
resistance emerges in virtually all patients within 9--18 months. The most common
resistance mechanism is the \textbf{C797S substitution} at the covalent binding
site (Cys797 $\to$ Ser797), which abrogates the irreversible covalent bond
formation required for osimertinib activity. When C797S occurs in \textit{cis}
with T790M (on the same allele), the resulting \textbf{triple-mutant EGFR} forms
(Ex19del/T790M/C797S and L858R/T790M/C797S) are \textit{resistant to all three
generations} of currently approved EGFR TKIs. Additional resistance mechanisms
include MET amplification, HER2 amplification, RAS/RAF mutations, and SCLC
histological transformation \parencite{leonetti2019resistance}. As of 2026, there
is no approved therapy for triple-mutant EGFR, and it remains a major long-term unmet
need in NSCLC treatment. It is important, however, to state the scope of this thesis
precisely at the outset: because no large-scale bioactivity dataset yet exists for the
triple mutant, the present work develops and validates its computational framework on
the abundant \emph{aggregate} EGFR inhibition data available in ChEMBL
(CHEMBL203)---which is dominated by wild-type and single-mutant measurements. The
framework therefore targets \emph{general} EGFR inhibitor identification; extending it
to mutant-specific prediction, including the triple mutant, is a direction for future
work rather than a claim of this study.

\subsection{Computational Drug Discovery and Graph Neural Networks}

Traditional drug discovery is a slow, costly, and attrition-prone process.
The average time from target identification to regulatory approval is 10--15 years,
and the average cost per approved drug has been estimated at approximately USD
2.0--2.6 billion \parencite{schneider2020rethinking}. The high attrition rate
($\sim$90\% of compounds entering Phase I clinical trials fail) is driven primarily
by inadequate efficacy, unexpected toxicity, and poor pharmacokinetic properties,
all of which are increasingly addressable through computational prediction.

\textbf{Computer-Aided Drug Design (CADD)} encompasses a broad spectrum of
computational techniques including molecular docking, molecular dynamics simulation,
free energy perturbation, QSAR modelling, and virtual screening. Among these,
machine learning-based virtual screening has emerged as the most scalable approach
for the rapid prioritisation of large compound libraries \parencite{walters2020generative}.

\textbf{Graph Neural Networks (GNNs)} represent a particularly powerful class of
machine learning architectures for molecular property prediction \parencite{gilmer2017neural}.
The fundamental insight is that molecules can be represented as mathematical graphs
in which atoms correspond to nodes carrying feature vectors and chemical bonds
correspond to edges carrying their own feature vectors, as illustrated in
Figure~\ref{fig:intro_mol2graph}. GNNs learn molecular representations through
iterative message-passing: at each layer, every atom aggregates feature information
from its bonded neighbours and updates its own hidden representation accordingly.
After several such message-passing layers, a readout step aggregates all the atom
representations into a single fixed-length molecular embedding that is passed to a
classifier or regression head.

\begin{figure}[htbp]
\centering
\includegraphics[width=0.95\textwidth]{molecule_to_graph.pdf}
\caption{Representation of a molecule as an attributed graph for graph neural
         network learning. The 2D structure of the EGFR inhibitor gefitinib (a) is
         converted into a molecular graph (b) in which each atom becomes a node and
         each bond an edge. Aromatic atoms are shown in green and non-aromatic atoms
         in blue; each atom and bond is encoded by a fixed-length feature vector.}
\label{fig:intro_mol2graph}
\end{figure}

This approach offers several advantages over traditional fingerprint-based molecular
representations. First, GNNs learn task-specific molecular representations
\textit{end-to-end} from bioactivity data, rather than relying on pre-defined
structural patterns encoded in fingerprints. Second, by incorporating edge feature
vectors, GNNs can natively encode bond-level information (bond order, conjugation,
stereochemistry) that is either discarded or only coarsely approximated in
circular fingerprints. Third, deep GNN architectures can capture long-range
atomic interactions spanning multiple bonds, enabling them to model higher-order
pharmacophoric relationships that influence molecular activity \parencite{jiang2021gnn}.

Despite these advantages, the translation of GNNs from academic benchmarks to
practical drug discovery applications faces challenges, including the need for
rigorous evaluation protocols, chemical interpretability, and extension to
prospective screening pipelines---all of which are directly addressed by this thesis.

\subsection{Nepali Medicinal Plants as a Natural-Product Resource}

Nepal's Himalayan biodiversity spans roughly 6,500 flowering plant species, of which
more than 700 have documented ethnopharmacological use in Ayurvedic and Tibetan
medicine. Many are inexpensive, culturally familiar, and locally accessible, making
them an attractive yet under-explored source of natural-product leads for
resource-limited settings. The prospective screening dataset used in this thesis is
drawn from \textbf{seven traditional Nepali medicinal plant species}---\textit{Swertia
chirayita}, \textit{Rhododendron arboreum}, \textit{Nardostachys jatamansi},
\textit{Terminalia chebula}, \textit{Berberis aristata}, \textit{Ocimum sanctum}, and
\textit{Tinospora cordifolia}---selected for their documented anti-inflammatory and
antiproliferative use and their coverage in the \textbf{COCONUT~2.0} open
natural-product database \parencite{sorokina2021coconut}. After RDKit standardisation
this yields a library of \textbf{535 unique natural-product compounds} (ranging from 9
compounds for \textit{Ocimum sanctum} to 203 for \textit{Nardostachys jatamansi}) that
is screened for EGFR inhibitor activity in Chapter~4. Screening this locally accessible
chemical space directly connects computational drug discovery to Nepal's biodiversity
and offers an affordable route to candidate anti-NSCLC compounds.

% ----------------------------------------------------------------
\section{Problem Statement}
% ----------------------------------------------------------------

Despite the transformative impact of EGFR-targeted therapy and the promise of
Graph Neural Networks for drug discovery, critical and interrelated gaps persist
in both clinical practice and computational methodology.

\subsection{Clinical and Translational Need}

Acquired resistance to EGFR tyrosine kinase inhibitors---most notably the
\textbf{C797S mutation}, which renders third-generation covalent inhibitors
ineffective and for which no fourth-generation drug is yet approved---sustains a
continuing need to discover structurally novel, non-covalent EGFR inhibitor
chemotypes \parencite{rotow2023egfr}. This need is most acute in South and East Asia,
where EGFR-mutant NSCLC is highly prevalent yet novel targeted therapies remain
expensive and computational drug-discovery infrastructure is scarce. There is
therefore a clear demand for open-source, CPU-executable screening pipelines that can
prioritise affordable, regionally accessible natural-product leads---the translational
problem this thesis addresses.

\subsection{Computational Problem: Methodological Gaps in GNN-Based EGFR Inhibitor Prediction}

A review of the existing computational literature reveals three specific methodological
deficiencies that this thesis is designed to address:

\begin{enumerate}
    \item \textbf{Evaluation bias from random splitting.} The majority of published
    GNN models for molecular property prediction, including those applied to EGFR
    inhibitor classification, employ random train/test splitting. This approach allows
    structurally similar or near-identical molecular scaffolds to appear simultaneously
    in the training and test sets. As a result, the reported performance metrics
    substantially overestimate a model's ability to generalise to genuinely novel
    chemical matter---precisely the task required in prospective virtual screening.
    A model that has ``memorised'' a scaffold from the training set cannot be trusted
    to correctly classify a structurally unprecedented compound in a real screening
    campaign.

    \item \textbf{Limited molecular featurisation.} Many published GNN models for
    EGFR or related kinase targets use minimal node feature sets (e.g., atomic number
    only, or a small set of atom type flags) and ignore edge (bond) features entirely.
    This discards potentially discriminative chemical information including bond order,
    conjugation, aromaticity, and stereochemistry---all of which influence molecular
    recognition in the EGFR binding pocket.

    \item \textbf{Absence of prospective virtual screening on natural product libraries.}
    Most academic GNN studies in the EGFR domain report classification performance on
    held-out test sets but do not extend their trained models to the most practically
    relevant application: scoring and ranking an external library of unseen compounds
    to identify novel lead candidates---especially compounds derived from traditional
    medicinal plants that may offer affordable, regionally accessible therapeutic options
    \parencite{malik2025deepegfr}.
\end{enumerate}

\noindent This thesis addresses these gaps through its two objectives. First, it
\emph{builds graph neural network classifiers}---with rich 29-dimensional atom and
12-dimensional bond featurisation, evaluated under a strict Bemis-Murcko scaffold
split---and \emph{benchmarks them against traditional machine-learning baselines}
(Random Forest and an MLP on ECFP4 fingerprints), directly testing whether the added
complexity of GNNs is justified over simpler, well-established models (gaps~1--2).
Second, it \emph{deploys the best-performing models to screen} a Nepali medicinal-plant
compound library for drug-like EGFR inhibitor leads (gap~3). Because large-scale
triple-mutant bioactivity data are not yet available, the models are trained on
aggregate EGFR data, and the natural-product scaffolds they surface are intended as
candidate starting points for future experimental validation rather than as validated
triple-mutant inhibitors.

% ----------------------------------------------------------------
\section{Objectives}
% ----------------------------------------------------------------

\subsection{General Objective}

To build and benchmark a Graph Neural Network framework for EGFR inhibitor
classification against traditional machine-learning baselines, and to deploy the
best-performing models for prospective virtual screening of Nepali medicinal plant
compounds as candidate anti-NSCLC leads.

\subsection{Specific Objectives}

\begin{enumerate}
    \item To build graph neural network models that classify EGFR inhibitors and
          compare them with traditional machine-learning models.

    \item To use the best models to screen Nepali medicinal plant compounds and
          identify drug-like EGFR inhibitor leads.
\end{enumerate}

% ----------------------------------------------------------------
\section{Rationale}
% ----------------------------------------------------------------

\textbf{Clinical rationale.}
Acquired resistance to EGFR TKIs---and ultimately the triple-mutant EGFR that follows
osimertinib failure---continues to drive the search for structurally novel inhibitor
chemotypes, yet bringing new therapies through conventional drug discovery takes years
to decades. Computational virtual screening can accelerate the early stage of this
search by pre-selecting the most promising candidates for experimental testing,
reducing the number of compounds that must be synthesised and assayed. In Nepal and
South Asia, where EGFR mutations are prevalent yet targeted therapies remain expensive,
this further motivates screening locally accessible, low-cost natural-product chemical
space rather than costly synthetic libraries \parencite{rotow2023egfr}.

\textbf{Scientific rationale.}
Graph Neural Networks have demonstrated strong performance across molecular
machine learning benchmarks \parencite{hu2020strategies}, yet their application
to EGFR inhibitor discovery has been constrained by the three methodological
shortcomings enumerated in Section~1.2. Addressing all three simultaneously within a
single, reproducible pipeline---scaffold splitting, rich featurisation, systematic
multi-architecture comparison, and prospective Nepali plant compound screening---is
what motivates the design of this study.

\textbf{Data science rationale.}
This project integrates multiple core data science competencies within a single
applied research context: large-scale database querying and API integration (ChEMBL),
cheminformatics feature engineering (RDKit), deep learning architecture design and
training (PyTorch, PyTorch Geometric), automated hyperparameter search (Optuna),
prospective natural compound screening (COCONUT 2.0), and scientific visualisation
(Matplotlib, Seaborn, RDKit 2D depictions). The project therefore constitutes a comprehensive demonstration
of advanced data science methodology applied to a high-impact biomedical problem.

\textbf{Practical rationale.}
The pipeline was deliberately built from exclusively open-source Python libraries
(RDKit, PyTorch Geometric, Optuna, scikit-learn) and public databases (ChEMBL,
COCONUT~2.0), running on commodity CPU hardware, so that every reported result can be
reproduced without licensed software or specialised computing resources.

% ----------------------------------------------------------------
\section{Research Hypotheses}
% ----------------------------------------------------------------

The following two testable research hypotheses were formulated prior to
experimentation and are evaluated against the results presented in Chapter 4:

\begin{enumerate}
    \item \textbf{$H_1$ (Performance Hypothesis):} Graph Neural Network architectures
    incorporating comprehensive multi-dimensional atomic node features and bond-level
    edge features, when optimised via Bayesian hyperparameter search, will achieve
    measurably superior scaffold test-set AUROC compared to traditional machine learning
    baselines (Random Forest and MLP on 2048-bit ECFP4 fingerprints)
    under a strict Bemis-Murcko scaffold split on the ChEMBL EGFR bioactivity dataset.

    \item \textbf{$H_2$ (Virtual Screening Hypothesis):} The virtual screening
    pipeline, employing the study's two best-performing models (the best graph network
    and the best overall classifier) to score compounds from seven traditional Nepali
    medicinal plant species retrieved from COCONUT 2.0, will identify---under
    \emph{each} of these models---at least one fully Lipinski-compliant
    natural-product-derived lead compound with a predicted activity probability
    $\hat{p} \geq 0.80$ and a Quantitative Estimate of Drug-likeness (QED) score
    $\geq 0.40$, contributing a viable experimental lead candidate for affordable NSCLC
    therapy development in the South Asian context.
\end{enumerate}

% ----------------------------------------------------------------
\section{Significance of the Study}
% ----------------------------------------------------------------

This research makes several original and complementary contributions to the
intersection of computational data science and molecular drug discovery:

\textbf{Methodological contribution.} Although systematic comparisons of GNN
architectures on molecular benchmarks are well established
\parencite{jiang2021gnn}, this thesis applies such a comparison, to our knowledge for
the first time, \emph{specifically to EGFR inhibitor classification}: it evaluates
five GNN architectures (GCN, GAT, GIN, AttentiveFP, GraphSAGE) against traditional
fingerprint baselines (Random Forest and MLP) under a rigorous Bemis-Murcko scaffold
split, using a rich 29-dimensional atom and 12-dimensional bond featurisation and
Bayesian hyperparameter optimisation (Optuna) for reproducible model selection. The
methodological novelty lies in this EGFR-specific, scaffold-aware, baseline-anchored
evaluation---coupled directly to the prospective natural-product screen below---rather
than in the architecture comparison in isolation.

\textbf{Natural product screening contribution.} By curating bioactive compounds
from seven traditional Nepali medicinal plant species (COCONUT 2.0) and applying the
study's best-performing models (the best graph network and the best overall
classifier) for prospective EGFR inhibition prediction, this thesis delivers the
first computational evaluation of Nepali natural products as potential EGFR TKI
candidates---directly relevant to the South Asian oncology community where EGFR
mutations are prevalent and novel therapies remain expensive.

\textbf{Translational contribution.} By extending the best-performing models into a
complete multi-model prospective virtual screening pipeline applied to
Nepali medicinal plant compounds and generating a ranked list of
Lipinski-compliant natural-product-derived
lead candidates with computed ADMET-proxy properties, the thesis delivers a concrete,
actionable output---a prioritised shortlist of natural compounds for future
experimental validation---rather than a purely academic performance comparison.

\textbf{Regional and social contribution.} The elevated prevalence of EGFR-activating
mutations in East and South Asian NSCLC patients, combined with the rising cancer
burden and limited access to novel therapies in Nepal, makes this research
particularly relevant to the South Asian academic and clinical community. The
entirely open-source, CPU-executable implementation ensures that the methodology
can be adopted and extended by researchers in resource-limited settings without
access to expensive commercial software or high-performance computing clusters.

% ----------------------------------------------------------------
\section{Limitations}
% ----------------------------------------------------------------

The following limitations of the current study are acknowledged and should be
considered when interpreting the results:

\begin{itemize}
    \item \textbf{No experimental validation.} All compound classifications and
    virtual screening predictions are purely computational. No \textit{in vitro}
    biochemical assay, cell-based proliferation assay, or \textit{in vivo}
    pharmacological experiment was conducted within the scope of this thesis.
    The predicted lead compounds remain computational hypotheses requiring
    experimental confirmation before any biological conclusions can be drawn.

    \item \textbf{Bioassay data heterogeneity.} Bioactivity data sourced from
    ChEMBL aggregates IC$_{50}$ measurements originating from hundreds of
    independent laboratories employing diverse assay formats, cell lines, and
    experimental conditions. This inherent variability introduces label noise into
    the training dataset. While pIC$_{50}$ averaging and strict data filtering
    mitigate this to some extent, residual heterogeneity may impose a ceiling on
    achievable test-set performance under scaffold splitting.

    \item \textbf{Wild-type EGFR specificity.} No dedicated, large-scale bioactivity
    dataset exists for triple-mutant EGFR (Ex19del/T790M/C797S) in ChEMBL.
    Consequently, the model was trained on the aggregate EGFR bioactivity data of
    CHEMBL203, which is dominated by wild-type and single-mutant measurements. While
    the prospective screening of structurally novel Nepali plant scaffolds is designed
    to surface chemotypes distinct from existing covalent inhibitors, the model cannot
    explicitly predict activity against the C797S mutant form.

    \item \textbf{Incomplete database coverage of Nepali medicinal flora.} The
    virtual screening library is limited to compounds recorded for the seven target
    species in the COCONUT~2.0 database. Coverage is uneven across species and many
    traditional Nepali plant compounds remain uncharacterised and absent from public
    cheminformatics databases, so the screened library is not exhaustive of the
    phytochemical diversity of these plants.

    \item \textbf{Model-sensitivity of the virtual screen.} The classifiers were
    trained on synthetic kinase inhibitors but applied to structurally dissimilar
    natural products, and the resulting hit list depends strongly on the model. As
    reported in Chapter~4, the number of predicted actives and the identity of the top
    leads differ markedly between the tuned AttentiveFP, the best graph network (GCN),
    and the best classifier (Random Forest). No single model can be treated as
    authoritative; the study mitigates this by deploying the two best models and
    prioritising cross-model consensus leads, but only experimental assay can resolve
    the residual uncertainty.

    \item \textbf{Out-of-distribution prediction reliability.} On chemical space far
    from the training distribution, the weakest model (AttentiveFP) assigned
    near-saturated probabilities to large polyphenolic tannins that fail drug-likeness
    filters, and even the best models still ranked some large, non-drug-like
    polyphenols and triterpenoids highly---indicating that such predictions must be
    interpreted with caution and coupled with physicochemical filtering.

    \item \textbf{Two-dimensional molecular representation.} The molecular graph
    encodes the constitutional (2D) topology of the molecule: atom identity and
    bonding connectivity. Three-dimensional conformational effects, induced-fit
    protein-ligand binding geometry, explicit solvent interactions, and
    protein-mediated allosteric effects are not represented in the input features.

    \item \textbf{Binary classification scope.} The study frames EGFR inhibition
    as a binary classification problem. This formulation discards the continuous
    nature of the pIC$_{50}$ scale and cannot distinguish between a moderately
    potent compound ($\text{pIC}_{50} = 6.1$) and a highly potent compound
    ($\text{pIC}_{50} = 9.5$). Continuous regression modelling of pIC$_{50}$
    values is beyond the current scope but is recommended for future work.

    \item \textbf{ADMET-proxy limitations.} The ADMET profiling in this study
    relies solely on RDKit-computed physicochemical descriptors (MW, LogP, HBD,
    HBA, TPSA, QED) as proxies for oral bioavailability. These filters do not
    capture metabolic stability, cytochrome P450 inhibition, plasma protein binding,
    blood-brain barrier permeability, hERG channel toxicity, or other critical
    ADMET parameters that determine clinical development success.
\end{itemize}

% ----------------------------------------------------------------
\section{Organisation of the Thesis}
% ----------------------------------------------------------------

The remainder of this thesis is organised as follows. \textbf{Chapter 2}
provides a comprehensive review of the relevant literature, covering the biology
of EGFR and NSCLC, the clinical development and resistance mechanisms of EGFR
TKIs, public bioactivity databases, molecular representation methods, GNN
architectures for drug discovery, traditional Nepali medicinal plants as a source
of EGFR inhibitor candidates, and virtual screening approaches. \textbf{Chapter 3} describes the complete materials and methods of the
five-phase computational pipeline, including data acquisition, molecular graph
construction, model architectures, training protocol, hyperparameter optimisation,
and the Nepali medicinal plant virtual screening workflow. \textbf{Chapter 4}
presents all experimental results with tables and analytical interpretation,
covering dataset characterisation, model benchmarking, hyperparameter tuning
outcomes, and Nepali plant virtual screening lead prioritisation.
\textbf{Chapter 5} draws the principal conclusions of the thesis and recommends
directions for future experimental and computational work.